\documentclass[a4paper]{arrowhead}

\usepackage[yyyymmdd]{datetime}
\usepackage{etoolbox}
\usepackage[utf8]{inputenc}
\usepackage[T1]{fontenc}
\usepackage{multirow}
\usepackage{longtable}
\usepackage{tikz}
\usepackage{adjustbox}
\AtBeginEnvironment{longtable}{\small}
\BeforeBeginEnvironment{tikzpicture}{\begin{adjustbox}{max width=\textwidth}}
\AfterEndEnvironment{tikzpicture}{\end{adjustbox}}
\usetikzlibrary{positioning,shapes.multipart,arrows.meta}

\renewcommand{\dateseparator}{-}
\setlength{\emergencystretch}{3em}
\lstset{language=HTTP,breakatwhitespace=false,columns=fullflexible}

%% Special references
\newcommand{\mref}[2][]{{\textcolor{ArrowheadPurple}{\hyperref[sec:model:#2]{\ifstrempty{#1}{#2}{#1}}}}}
\newcommand{\oref}[1]{{\textcolor{ArrowheadBlue}{\hyperref[sec:operations:#1]{#1}}}}
\newcommand{\pdef}[1]{{\textcolor{ArrowheadGrey}{#1\label{sec:model:primitives:#1}\label{sec:model:primitives:#1s}\label{sec:model:primitives:#1es}}}}
\newcommand{\pref}[1]{{\textcolor{ArrowheadGrey}{\hyperref[sec:model:primitives:#1]{#1}}}}

% as template_sd.tex; the input and output arguments are passed already formatted (\mref or \pref)
\newrobustcmd\osubsection[3]{
  \addtocounter{subsection}{1}
  \addcontentsline{toc}{subsection}{\protect\numberline{\thesubsection}operation \textcolor{ArrowheadBlue}{#1}}
  \subsection*{
    \thesubsection\quad
    operation
    \textcolor{ArrowheadBlue}{#1}
    (\notblank{#2}{#2}{})
    \notblank{#3}{: #3}{}
  }
  \label{sec:operations:#1}
}
\newrobustcmd\msubsection[2]{
  \addtocounter{subsection}{1}
  \addcontentsline{toc}{subsection}{\protect\numberline{\thesubsection}#1 \textcolor{ArrowheadPurple}{#2}}
  \subsection*{\thesubsection\quad#1 \textcolor{ArrowheadPurple}{#2}}
  \label{sec:model:#2} \label{sec:model:#2es}
}
%%

\usepackage{xurl}
\begin{document}
\sloppy

%% Arrowhead Document Properties
\ArrowheadTitle{authorization}
\ArrowheadServiceID{authorization}
\ArrowheadType{Service Description}
\ArrowheadTypeShort{SD}
\ArrowheadVersion{5.3.0}
\ArrowheadDate{2026-11-09}
\ArrowheadAuthor{Jerker Delsing}
\ArrowheadStatus{REVIEW DRAFT}
\ArrowheadContact{jerker.delsing@ltu.se}
\ArrowheadFooter{\href{www.arrowhead.eu}{www.arrowhead.eu}}
\ArrowheadSetup
%%

%% Front Page
\begin{center}
  \vspace*{1cm}
  \huge{\arrowtitle}

  \vspace*{0.2cm}
  \LARGE{\arrowtype}
  \vspace*{1cm}

  \Large{Service ID: \textit{"\arrowid"}}

  \vspace*{\fill}

  % Front Page Image

  \vspace*{1cm}
  \vspace*{\fill}

  % Front Page Abstract
  \begin{abstract}
    This document describes the \texttt{authorization} service, version 1.0.0, produced by the ConsumerAuthorization core system of Arrowhead (documentation edition 5.3.0).
  \end{abstract}

  \vspace*{1cm}
\end{center}

\newpage
%%

%% Table of Contents
\tableofcontents
\newpage
%%

\section{Documentation edition and implementation scope}
This is the 5.3.0 documentation edition, dated 2026-11-09. It carries forward the system and service contracts from the retained 5.2.2 documentation and identifies the current primary implementation sources. The former \path|documentation/v5.2.2| folder is retained unchanged.

Document version 5.3.0 is distinct from executable, dependency, registered system/service and wire-contract versions. Values still defined by the source (including 5.2.2 artefact versions, 1.0.0 service versions and \path|SysMonitor-AMQP-IDD-5.2.2| payload identifiers) are preserved. This documentation update does not change runtime behaviour or establish a new implementation release. Historical verification statements describe the tests recorded in the inherited baseline; this revision validates document compilation and references.

\section{Overview}
\label{sec:overview}

This document describes the \texttt{authorization} service, which provides provider-level authorization policies: a provider declares which consumers may use its service (or event type), and any system can ask whether a consumer is authorized.
It is produced by the ConsumerAuthorization system \cite{ca-sysd} in service version 1.0.0.

The rest of this document is organized as follows.
In Section \ref{sec:operations}, we describe the abstract message operations provided by the service.
In Section \ref{sec:model}, we end the document by presenting the data types used by the mentioned operations.
The concrete interface is specified in the HTTP/TLS/JSON IDD \cite{ca-authorization-http-idd}. The MQTT binding is specified in the MQTT/TLS/JSON IDD \cite{ca-authorization-mqtt-idd}.

\subsection{Significant Prior Art}

Successor of the Arrowhead 4.x Authorization \texttt{intracloud} rules. In 5.x a policy has a default rule and optional per-operation (scoped) rules, each \texttt{ALL}, \texttt{WHITELIST}, \texttt{BLACKLIST} or \texttt{SYS\_\allowbreak{}METADATA}.

\subsection{How This Service Is Meant to Be Used}

A provider calls \texttt{grant} to publish the policy of one of its services and \texttt{revoke} to remove it; \texttt{lookup} lists its own policies. Orchestrators and providers call \texttt{verify} to evaluate a consumer request against the policies.

\subsection{Important Delimitations and Dependencies}
\label{sec:delimitations}

A provider can only set policies for itself (level \texttt{PROVIDER}). Management-level policies (level \texttt{MGMT}), set through \texttt{authorization\allowbreak{}Management}, take precedence: provider policies are only evaluated when no management policy applies. The service does not enforce anything; enforcement is done by the orchestrator (which only returns authorized providers) and by providers that check tokens.

\newpage

\section{Service Interface}
\label{sec:operations}

This section describes the interfaces to the \texttt{authorization} service.
In particular, each subsection names an abstract operation, an input type and an output type, in that order.
The input type is named inside parentheses, while the output type is preceded by a colon.
Input and output types are only denoted when accepted or returned, respectively, by the interface in question.

All abstract data types named in this section are defined in Section \ref{sec:model}.

\begin{figure}[ht!]
  \centering
  \begin{tikzpicture}[font=\small]
    \node[draw=ArrowheadBlue, very thick, rectangle split, rectangle split parts=2, align=center,
          rectangle split part fill={ArrowheadBlue!12,white}, text width=9cm, inner sep=6pt] (b)
      {\guillemotleft{}Service\guillemotright{}\\\textbf{authorization}
        \nodepart{second} \texttt{grant}\\ \texttt{revoke}\\ \texttt{lookup}\\ \texttt{verify}};
  \end{tikzpicture}
  \caption{SysML-style block of the \texttt{authorization} service and its operations.}
  \label{fig:SD}
\end{figure}

The following interface operations are available.

\osubsection{grant}{\mref[Authorization\allowbreak{}Grant\allowbreak{}Request]{AuthorizationGrantRequest}}{\mref[Authorization\allowbreak{}Policy\allowbreak{}Response]{AuthorizationPolicyResponse}}

Creates or replaces the provider-level policy of the requesting provider for one target. Returns 201 for a new policy, 200 if an identical one exists.

\osubsection{revoke}{\pref{Identifier}}{}

Removes a provider-level policy of the requester by instance identifier. Returns 200 if removed, 204 if it did not exist.

\osubsection{lookup}{\mref[Authorization\allowbreak{}Lookup\allowbreak{}Request]{AuthorizationLookupRequest}}{\mref[Authorization\allowbreak{}Policy\allowbreak{}List\allowbreak{}Response]{AuthorizationPolicyListResponse}}

Returns the requester's provider-level policies matching the filter.

\osubsection{verify}{\mref[Authorization\allowbreak{}Verify\allowbreak{}Request]{AuthorizationVerifyRequest}}{\pref{Boolean}}

Evaluates whether the consumer may use the target (and operation) of the provider: management-level policies first, provider-level policies if none applies; no applicable policy means denied. Returns a boolean.

\newpage

\section{Information Model}
\label{sec:model}

Here, all data objects that can be part of the \texttt{authorization} service are listed in alphabetic order.
Note that each subsection, which describes one type of object, begins with the \textit{struct} keyword, which is used to denote a collection of named fields, each with its own data type.
As a complement to the explicitly defined types in this section, there is also a list of implicit primitive types in Section \ref{sec:model:primitives}, which are used to represent things like hashes and identifiers.

\msubsection{struct}{AuthorizationGrantRequest}

A provider-level authorization policy that the provider sets for its own service.


\noindent\begin{tabularx}{\textwidth}{| >{\raggedright\arraybackslash}p{3.9cm} | >{\raggedright\arraybackslash\small}p{5.3cm} | X |} \hline
\rowcolor{gray!33} Field & Type & Description \\ \hline
cloud & \pref{Name} & Cloud identifier of the consumer cloud the policy applies to; \texttt{LOCAL} for the own cloud. \\ \hline
target\allowbreak{}Type & \pref{String} & Target type: \texttt{SERVICE\_\allowbreak{}DEF} or \texttt{EVENT\_\allowbreak{}TYPE}. \\ \hline
target & \pref{Name} & Service definition (or event type) the policy protects. \\ \hline
description & \pref{String} & Free-text description. \\ \hline
default\allowbreak{}Policy & \mref[Authorization\allowbreak{}Policy\allowbreak{}Request]{AuthorizationPolicyRequest} & Policy that applies to all operations without a scoped policy. \\ \hline
scoped\allowbreak{}Policies & \pref{Map}$<$\allowbreak{}\pref{String}, \mref[Authorization\allowbreak{}Policy\allowbreak{}Request]{AuthorizationPolicyRequest}$>$ & Policies for individual operations, keyed by operation name. \\ \hline
\end{tabularx}



\msubsection{struct}{AuthorizationLookupRequest}

Filter for the requester's own provider-level policies.


\noindent\begin{tabularx}{\textwidth}{| >{\raggedright\arraybackslash}p{3.9cm} | >{\raggedright\arraybackslash\small}p{5.3cm} | X |} \hline
\rowcolor{gray!33} Field & Type & Description \\ \hline
instance\allowbreak{}Ids & \pref{List}$<$\allowbreak{}\pref{Identifier}$>$ & Policy instance identifiers to match. \\ \hline
cloud\allowbreak{}Identifiers & \pref{List}$<$\allowbreak{}\pref{String}$>$ & Cloud identifiers to match. \\ \hline
target\allowbreak{}Names & \pref{List}$<$\allowbreak{}\pref{String}$>$ & Target names to match. \\ \hline
target\allowbreak{}Type & \pref{String} & Target type: \texttt{SERVICE\_\allowbreak{}DEF} or \texttt{EVENT\_\allowbreak{}TYPE}. \\ \hline
\end{tabularx}



\msubsection{struct}{AuthorizationPolicy}

One policy rule.


\noindent\begin{tabularx}{\textwidth}{| >{\raggedright\arraybackslash}p{3.9cm} | >{\raggedright\arraybackslash\small}p{5.3cm} | X |} \hline
\rowcolor{gray!33} Field & Type & Description \\ \hline
policy\allowbreak{}Type & \pref{AuthorizationPolicyType} & Policy type: \texttt{ALL}, \texttt{WHITELIST}, \texttt{BLACKLIST} or \texttt{SYS\_\allowbreak{}METADATA}. \\ \hline
policy\allowbreak{}List & \pref{List}$<$\allowbreak{}\pref{String}$>$ & System names of the white- or blacklist. \\ \hline
policy\allowbreak{}Metadata\allowbreak{}Requirement & \mref{MetadataRequirement} & Metadata the consumer system must have (\texttt{SYS\_\allowbreak{}METADATA}). \\ \hline
\end{tabularx}



\msubsection{struct}{AuthorizationPolicyListResponse}

List of authorization policies.


\noindent\begin{tabularx}{\textwidth}{| >{\raggedright\arraybackslash}p{3.9cm} | >{\raggedright\arraybackslash\small}p{5.3cm} | X |} \hline
\rowcolor{gray!33} Field & Type & Description \\ \hline
entries & \pref{List}$<$\allowbreak{}\mref[Authorization\allowbreak{}Policy\allowbreak{}Response]{AuthorizationPolicyResponse}$>$ & Matching records of the current page. \\ \hline
count & \pref{Number} & Total number of matching records (all pages). \\ \hline
\end{tabularx}



\msubsection{struct}{AuthorizationPolicyRequest}

One policy rule.


\noindent\begin{tabularx}{\textwidth}{| >{\raggedright\arraybackslash}p{3.9cm} | >{\raggedright\arraybackslash\small}p{5.3cm} | X |} \hline
\rowcolor{gray!33} Field & Type & Description \\ \hline
policy\allowbreak{}Type & \pref{String} & Policy type: \texttt{ALL}, \texttt{WHITELIST}, \texttt{BLACKLIST} or \texttt{SYS\_\allowbreak{}METADATA}. \\ \hline
policy\allowbreak{}List & \pref{List}$<$\allowbreak{}\pref{String}$>$ & System names of the white- or blacklist. \\ \hline
policy\allowbreak{}Metadata\allowbreak{}Requirement & \mref{MetadataRequirement} & Metadata the consumer system must have (\texttt{SYS\_\allowbreak{}METADATA}). \\ \hline
\end{tabularx}



\msubsection{struct}{AuthorizationPolicyResponse}

A stored authorization policy.


\noindent\begin{tabularx}{\textwidth}{| >{\raggedright\arraybackslash}p{3.9cm} | >{\raggedright\arraybackslash\small}p{5.3cm} | X |} \hline
\rowcolor{gray!33} Field & Type & Description \\ \hline
instance\allowbreak{}Id & \pref{Identifier} & Policy instance identifier \texttt{\textless{}MGMT or PR\textgreater{}\textbar{}\allowbreak{}\textless{}cloud\textgreater{}\textbar{}\allowbreak{}\textless{}provider\textgreater{}\textbar{}\allowbreak{}\textless{}target\allowbreak{}Type\textgreater{}\textbar{}\allowbreak{}\textless{}target\textgreater{}}. \\ \hline
level & \pref{AuthorizationLevel} & Policy level: \texttt{MGMT} (set by an operator) or \texttt{PROVIDER} (set by the provider itself). \\ \hline
cloud & \pref{Name} & Cloud identifier of the consumer cloud the policy applies to; \texttt{LOCAL} for the own cloud. \\ \hline
provider & \pref{Name} & Provider system the policy protects. \\ \hline
target\allowbreak{}Type & \pref{AuthorizationTargetType} & Target type: \texttt{SERVICE\_\allowbreak{}DEF} or \texttt{EVENT\_\allowbreak{}TYPE}. \\ \hline
target & \pref{Name} & Target of the policy or token: a service definition or event type. \\ \hline
description & \pref{String} & Free-text description. \\ \hline
default\allowbreak{}Policy & \mref{AuthorizationPolicy} & Policy that applies to all operations without a scoped policy. \\ \hline
scoped\allowbreak{}Policies & \pref{Map}$<$\allowbreak{}\pref{String}, \mref{AuthorizationPolicy}$>$ & Policies for individual operations, keyed by operation name. \\ \hline
created\allowbreak{}By & \pref{Name} & Name of the system that created the record. \\ \hline
created\allowbreak{}At & \pref{DateTime} & UTC time of creation. \\ \hline
\end{tabularx}



\msubsection{struct}{AuthorizationVerifyRequest}

One authorization question: may the consumer use the target (operation) of the provider?


\noindent\begin{tabularx}{\textwidth}{| >{\raggedright\arraybackslash}p{3.9cm} | >{\raggedright\arraybackslash\small}p{5.3cm} | X |} \hline
\rowcolor{gray!33} Field & Type & Description \\ \hline
provider & \pref{Name} & Provider system; for \texttt{verify} defaults to the requester. \\ \hline
consumer & \pref{Name} & Consumer system; for \texttt{verify} defaults to the requester. \\ \hline
cloud & \pref{Name} & Cloud of the consumer, \texttt{LOCAL} by default. \\ \hline
target\allowbreak{}Type & \pref{String} & Target type: \texttt{SERVICE\_\allowbreak{}DEF} or \texttt{EVENT\_\allowbreak{}TYPE}. \\ \hline
target & \pref{Name} & Target of the policy or token: a service definition or event type. \\ \hline
scope & \pref{Name} & Operation to check; empty checks the default policy. \\ \hline
\end{tabularx}



\msubsection{struct}{MetadataRequirement}

A JSON object of metadata requirements: each key is a metadata path (dot notation), each value the required value or an operator object such as \texttt{\{"op": "LESS\_\allowbreak{}THAN",\allowbreak{} "value": 10\}}.



\subsection{Primitives}
\label{sec:model:primitives}

Types and structures mentioned throughout this document that are assumed to be available to implementations of this service.
The concrete interpretations of each of these types and structures must be provided by any IDD document claiming to implement this service.

\noindent\begin{tabularx}{\textwidth}{| p{4.6cm} | X |} \hline
\rowcolor{gray!33} Type & Description \\ \hline
\pdef{AuthorizationLevel} & One of \texttt{MGMT}, \texttt{PROVIDER}. \\ \hline
\pdef{AuthorizationPolicyType} & One of \texttt{ALL}, \texttt{WHITELIST}, \texttt{BLACKLIST}, \texttt{SYS\_\allowbreak{}METADATA}. \\ \hline
\pdef{AuthorizationTargetType} & One of \texttt{SERVICE\_\allowbreak{}DEF}, \texttt{EVENT\_\allowbreak{}TYPE}. \\ \hline
\pdef{Boolean} & One out of \texttt{true} or \texttt{false}. \\ \hline
\pdef{DateTime} & Pinpoints a specific moment in time (UTC). \\ \hline
\pdef{Identifier} & A string that identifies a record, e.g. a UUID or a composite identifier. \\ \hline
\pdef{List} & An \textit{array} of a known number of items, each having type A. \\ \hline
\pdef{Map} & A set of key-value pairs. \\ \hline
\pdef{Name} & A string identifier that is intended to be both human and machine-readable. \\ \hline
\pdef{Number} & An integer number. \\ \hline
\pdef{String} & A chain of characters. \\ \hline
\end{tabularx}

\newpage

\bibliographystyle{IEEEtran}
\bibliography{bibliography}

\newpage

\section{Revision History}
\subsection{Amendments}

\noindent\begin{tabularx}{\textwidth}{| p{1cm} | p{3cm} | p{2cm} | X | p{4cm} |} \hline
\rowcolor{gray!33} No. & Date & Version & Subject of Amendments & Author \\ \hline

1 & 2026-10-22 & 5.2.2 & Document created from the Arrowhead LaTeX templates for the 5.2.2 Java Spring implementation, including the local extensions of this workspace. & Jerker Delsing \\ \hline
2 & 2026-11-09 & 5.3.0 & New documentation edition; primary-source references and Go/Rust implementation scope; 5.2.2 retained. & Documentation maintenance \\ \hline
\end{tabularx}

\subsection{Quality Assurance}

\noindent\begin{tabularx}{\textwidth}{| p{1cm} | p{3cm} | p{2cm} | X |} \hline
\rowcolor{gray!33} No. & Date & Version & Approved by \\ \hline

1 & & \arrowversion  &  \\ \hline

\end{tabularx}

\end{document}
